\documentclass[11pt]{article}
\usepackage[margin=1in]{geometry}
\usepackage{amsmath,booktabs,hyperref}
\title{Quarantined Query Slots: A Pessimistic Research Countermeasure to Lease-Expiry Over-Admission}
\author{AssistX fleet trace-index engineering research}
\date{October 8, 2026}
\begin{document}
\maketitle
\begin{abstract}
A Redis in-flight capacity lease that expires independently of query execution cannot provide a hard bound on concurrent database operations. Prior synthetic experiments in the AssistX trace investigation stack demonstrated an old query surviving lease expiry while a replacement was admitted. We investigate a conservative alternative: atomic non-expiring Redis hash slots released only after an asserted query-completion condition. In 109 focused synthetic Python tests and a disconnected disposable Redis 7 Lua probe, slot occupancy remained bounded across simulated elapsed time and exact-token releases. Two negative findings limit applicability: a crashed worker holds capacity indefinitely, and Redis state loss can admit a replacement while an old query is still active. The caller-supplied completion assertion is not independently authenticated. This research code is not installed in the production trace API and does not establish durable or physical concurrency fencing.
\end{abstract}
\section{Failure model}
For a capacity $C$, suppose query $q_1$ holds a Redis token $t_1$ with expiry $e_1$. If wall time exceeds $e_1$ without a verified Neo4j cancellation, automatic token eviction allows $q_2$ to acquire another slot, even while $q_1$ remains physically active. Then the Redis cardinality may be $C$ although the true number of active database requests exceeds $C$. A requested Neo4j transaction timeout cannot alone be treated as an end-to-end transport and cancellation acknowledgment. We distinguish logical Redis occupancy from physical execution.
\section{Prediction and method}
The pre-experiment prospectus proposed atomic Redis hash occupancy with no expiry or stale-token auto-reap. Let $H$ denote the token hash. Admission uses an atomic Redis Lua script: if $|H|\ge C$, deny; otherwise insert a random token and return admission. A separately verified query-completion event would authorize removal of exactly the original token. Since no such trusted independent completion event is available, the experiment's boolean assertion is not sufficient for production use.

We preregistered that elapsed time would not decrease $|H|$; a lost worker would retain its slot indefinitely; malicious/incorrect tokens would not remove another query's occupancy; and a Redis state reset would still violate the physical bound.
\section{Observed synthetic tests}
\begin{table}[h]\centering
\caption{Read-only operational authority; synthetic local evidence only.}
\begin{tabular}{ll}\toprule
Threat or criterion & Observation\\\midrule
Simulated time advancement beyond prior 30s TTL & No successor admitted\\
Concurrent cap=3 synthetic admission attempts & Exactly 3 tokens granted\\
Mismatched and stale release token & No other slot removed\\
Redis exception/malformed reply & Admission denied\\
Worker crash with unverified cancellation & Capacity remains quarantined\\
Redis state lost while query remains active & Over-admission reproduced\\
Scoped Python test suite & 118 passing\\
Disconnected Redis 7 Lua probe & Expected outcomes observed\\
Live Neo4j cancellation receipt and failover & Not evaluated\\\bottomrule
\end{tabular}
\end{table}
The real Redis probe ran solely in a new container with no network, published ports, bind mounts, production keys or durable data. It tested Redis \texttt{HLEN}, \texttt{HSET}, \texttt{HDEL}, no key TTL and a deliberate test-key deletion. The isolated fixture was stopped afterward.
\section{Safety interpretation}
Never-expiring capacity is a liveness sacrifice: a vanished query owner cannot automatically regain its slot. It improves one failure mode, expiry over-admission, only while Redis retains authoritative state. External Redis failover or split brain, accidental deletion, forged completion assertions, and unreconciled remote Neo4j execution remain open. A durable independently authenticated execution identity, loss-resistant fencing epoch and physical cancellation witness are needed before activation. Existing fleet log-custody obligations remain independent \cite{nist}.
\section{Conclusions}
This is an incremental research design, not a production rollout or an authorization to merge admission PRs. It avoids a known TTL-induced over-admission mechanism but does not prove a hard physical concurrency bound under Redis data loss or uncertain cancellation. All real fleet sources and services were left unchanged. This arXiv-style manuscript is an uncompiled, unsubmitted working draft.
\begin{thebibliography}{9}
\bibitem{nist} K. Kent and M. Souppaya, \emph{Guide to Computer Security Log Management}, NIST SP 800-92 (2006). \url{https://doi.org/10.6028/NIST.SP.800-92}.
\bibitem{redis} Redis, \emph{EVAL}, Redis command documentation. \url{https://redis.io/docs/latest/commands/eval/}.
\bibitem{neo} Neo4j, \emph{Transaction Configuration}, Neo4j Python driver documentation. \url{https://neo4j.com/docs/python-manual/current/transactions/}.
\end{thebibliography}
\end{document}
